
          %%%%% ~~~~~~~~~~~~~~~~~~~~ %%%%%

\section{Estimation}

          %%%%% ~~~~~~~~~~~~~~~~~~~~ %%%%%

\subsection{Area estimates for ecoprovinces $\times$ reference land cover types}
For each ecoprovince \,$p$\, and reference land cover type \,$r$,\,
we computed the estimate \,$A_{p}^{(r)}$\,
for the area of the reference land cover type \,$r$\, within the ecoprovince \,$p$\, as follows:
First, recall that stratification was done by ecoprovince $\times$ stratification land cover 
type. Let \,$k$\, be a stratification land cover type.
Then, \,$(p,k)$\, represents a stratum, and, based 
on the reference data, we computed the Horvitz-Thompson estimator
\,$T^{(r)}_{(p,k)}$\, for the total area of reference land cover type \,$r$\,
within the stratum \,$(p,k)$.\,
Next, the Horvitz-Thompson area estimator \,$T_{p}^{(r)}$\,
for the reference land cover type \,$r$\, within ecoprovince \,$p$\, is the sum of 
the \,$T^{(r)}_{(p,k)}$’s, summing over the stratification land cover types \,$k$.\,
The variance estimate of \,$T^{(r)}_{p}$\,
is the sum of the variance estimates of its summands,
since the summands of
\,$T^{(r)}_{p}$\,
are mutually independent.
Lastly, for each ecoprovince \,$p$,\, the estimates
\,$A_{p}^{(r)}$’s\,
for different reference land cover types \,$r$\, were obtained from the
\,$T_{p}^{(r)}$'s\,
by {\color{red}calibrating} the
\,$T_{p}^{(r)}$’s\,
to the known area of the ecoprovince \,$p$.\, 

\vskip 0.3cm
\noindent
Precisely, the Horvitz-Thompson estimator
\,$T^{(r)}_{p}$\,
can be given as follows:

%The Horvitz-Thompson estimator under stratified random sampling:
%\vskip -0.5cm
\begin{eqnarray*}
T^{(r)}_{p}
& := &
% 	\underset{i\,\in\,\mathcal{U}_{p}}{\sum}\;y^{(r)}_{i}
% \;\, \approx \;\;
% 	{\color{red}\widehat{A^{(r)}_{p}}^{\textnormal{HT}}}
% \;\, {\color{black}:=} \;\;
	{\color{black}\underset{i\,\in\mathcal{S}_{p}}{\sum}\,\dfrac{y^{(r)}_{i}}{\pi_{i}}}
\;\; = \;\;
	\overset{K}{\underset{k = 1}{\sum}}
	\left(
		\underset{i\,\in\mathcal{S}_{p,k}}{\sum}\dfrac{y^{(r)}_{i}}{\pi_{i}}
		\right),
	\quad
	\textnormal{noting that
		\,$T^{(r)}_{(p,k)} \,=\, \underset{i\,\in\mathcal{S}_{p,k}}{\sum}\dfrac{y^{(r)}_{i}}{\pi_{i}}$
		}
\\
& = &
	\overset{K}{\underset{k = 1}{\sum}}
	\left(
		\underset{i\,\in\mathcal{S}_{p,k}}{\sum}\,
		\dfrac{M_{p,k}}{m_{p,k}}\,y^{(r)}_{i}
		\right),
	\quad
	\textnormal{since \,$\pi_{i} \,=\, \dfrac{m_{p,k}}{M_{p,k}}$,\, for each \,$i \in \mathcal{S}_{p,k}$}
\\
& = &
	\overset{K}{\underset{k = 1}{\sum}}\;
	M_{p,k}
	\left(
		\dfrac{1}{m_{p,k}}
		\cdot\!
		\underset{i\,\in\mathcal{S}_{p,k}}{\sum}
		y^{(r)}_{i}
		\right)
\;\; = \;\;
	\overset{K}{\underset{k=1}{\sum}}\;
	M_{p,k}
	\cdot
	\widehat{\overline{y}}^{(r)}_{p,k}
\end{eqnarray*}

where
\begin{multicols}{2}
    \begin{center}
    \begin{minipage}{3.25in}
    \begin{eqnarray*}
    \mathcal{U}_{p}{\color{white}..}
    & := &
        \left\{\!\!\begin{array}{c}
        \textnormal{\small collection of all grid cells}
        \\
        \textnormal{\small in ecoprovince \,$p$}
        \end{array}\!\!\right\}
    \\
    k{\color{white}..}
    & = &
        1, 2, \ldots, K \;\;\textnormal{($K$ = \# map classes)}
    \\
    \mathcal{U}_{p,k}
    & := &
        \left\{\!\!\begin{array}{c}
        \textnormal{\small collection ({\color{black}stratum}) of all grid cells}
        \\
        \textnormal{\small of {\color{red}map class \,$k$}\, in ecoprovince \,$p$}
        \end{array}\!\!\right\}
    \\
    \mathcal{S}_{p,k}
    & := &
        \left\{\!\!\begin{array}{c}
        \textnormal{\small{\color{black}simple random sample} of grid cells}
        \\
        \textnormal{\small selected from \,$\mathcal{U}_{p,k}$}
        \end{array}\!\!\right\}
    \\
    \mathcal{S}_{p}
    & := &
        \overset{K}{\underset{k=1}{\bigsqcup}}\;\mathcal{S}_{p,k}
    \\
    \end{eqnarray*}
    \end{minipage}
    \end{center}
\columnbreak
    \begin{center}
    \begin{minipage}{3.25in}
    \begin{eqnarray*}
    M_{p,k}
    & := &
        \left\vert\;\,\overset{{\color{white}.}}{\mathcal{U}_{p,k}}\;\right\vert
    \;\, = \;\,
        \left\{\!\!\begin{array}{c}
        \textnormal{\small number of all grid cells}
        \\
        \textnormal{\small in stratum \,$\mathcal{U}_{p,k}$}
        \end{array}\!\!\right\}
    \\
    m_{p,k}
    & := &
        \left\vert\;\overset{{\color{white}.}}{\mathcal{S}_{p,k}}\;\right\vert
    \;\, = \;\,
        \left\{\!\!\begin{array}{c}
        \textnormal{\small number of all grid cells}
        \\
        \textnormal{\small in stratum \,$\mathcal{S}_{p,k}$}
        \end{array}\!\!\right\}
    \\
    y^{(r)}_{i}
    & := &
        \left\{\begin{array}{c}
        \textnormal{\small land area of {\color{red}reference class $r$}}
        \\
        \textnormal{\small in the $i^{\textnormal{\,th}}$ grid cell}
        \end{array}\!\right\}
    \\
    \pi_{i}
    & := &
        \left\{\begin{array}{c}
        \textnormal{\small selection probability}
        \\
        \textnormal{\small of the $i^{\textnormal{\,th}}$ grid cell}
        \end{array}\!\right\}
    \end{eqnarray*}
    \end{minipage}
    \end{center}
\end{multicols}
\noindent
And, the variance estimator of \,$T^{(r)}_{p}$\, can be given by:

%The Horvitz-Thompson estimator under stratified random sampling:
%\vskip -0.5cm
\begin{eqnarray*}
% \widehat{\Var}\!\left(\,\widehat{A^{(k)}_{r}}^{\textnormal{HT}}\,\right)
\widehat{\Var}\!\left(\;
	\overset{{\color{white}.}}{T^{(r)}_{p}}
	\,\right)
& = &
	\overset{K}{\underset{k=1}{\sum}}\;\,
	\widehat{\Var}\!\left(\,
		\overset{{\color{white}.}}{T^{(r)}_{p,k}}
		\,\right)
	% \left(\,1 \,-\, \dfrac{m_{p,k}}{M_{p,k}}\,\right)
	% \cdot
	% M_{p,k}^{2}
	% \cdot
	% \dfrac{s_{p,k}^{2}}{m_{p,k}}
\;\; = \;\;
	\overset{K}{\underset{k=1}{\sum}}
	\left(\,1 \,-\, \dfrac{m_{p,k}}{M_{p,k}}\,\right)
	\cdot
	M_{p,k}^{2}
	\cdot
	\dfrac{s_{p,k}^{2}}{m_{p,k}}
\end{eqnarray*}
where
\begin{equation*}
s_{p,k}^{2}
\; = \;
	\dfrac{1}{m_{p,k}\,-\,1}
	\cdot\!
	\underset{i\,\in\mathcal{S}_{p,k}}{\sum}\!
	\left(\,y^{(r)}_{i} \,-\, \overline{y}^{(k)}_{p,k}\,\right)^{2}
\quad\quad
\textnormal{and}
\quad\quad\;\;
\widehat{\overline{y}}^{(r)}_{p,k}
\; = \;
	\dfrac{1}{m_{p,k}}
	\cdot\!
	\underset{i\,\in\mathcal{S}_{p,k}}{\sum}
	y^{(r)}_{i}
\end{equation*}


          %%%%% ~~~~~~~~~~~~~~~~~~~~ %%%%%

\subsection{Area estimates for ecozones $\times$ reference land cover types}
Since the partition by ecoprovinces is a refinement of that by ecozones
(i.e., each ecozone is made up of its constituent ecoprovinces),
the area estimate for reference land cover type \,$r$\,
within ecozone \,$z$\, is the sum of the
\,$A_{p}^{(r)}$'s,\,
summing over the constituent ecoprovinces \,$p$\, within \,$z$.\,
The variance estimate for the Horvitz-Thompson area estimator for reference land 
cover type \,$r$\, within ecozone \,$z$\, is the sum of the variance estimates of its summands 
(summing over constituent ecoprovinces of \,$z$). 

          %%%%% ~~~~~~~~~~~~~~~~~~~~ %%%%%

\subsection{Area estimates for reference land cover types over all of Canada}
These were generated analogously as their ecoprovince and ecozone counterparts.
Since the land mass of Canada is partitioned into ecozones,
the area estimate for reference land cover type \,$r$\, over all of Canada
is the sum of the area estimates for reference land cover type \,$r$\,
over all the ecozones.
The variance estimate for the Horvitz-Thompson area estimator for 
reference land cover type \,$r$\, over all of Canada is the sum of the variance estimates of its 
summands (summing over all ecozones). 

          %%%%% ~~~~~~~~~~~~~~~~~~~~ %%%%%

\subsection{Area estimates for drainage regions $\times$ reference land cover types}
Recall that our stratification was done by ecoprovince and stratification land cover type.
On the other hand, the partition by ecoprovinces is not a refinement of that by drainage 
regions.
The preceding approach of generating ecozone-level estimates from ecoprovincelevel ones
(straightforward addition of estimates of constituent ecoprovinces) is not 
applicable.
Instead, the area estimate for a given drainage region \,$d$\, and reference land 
cover type \,$r$\, was computed as the weighted sum of the stratum-specific Horvitz-Thompson 
mean area estimates for the reference land cover type \,$r$\,, where the mean area estimate
for reference land cover type \,$r$\, within each stratum was weighted by the (non-random) 
``effective number of grid cells'' of that stratum within the drainage region \,$d$.\,
The ``effective number of grid cells'' of a given stratum within a given drainage region
is defined as the area of the intersection of the drainage region and the stratum
divided by the (constant) area of the grid cell (the grid cell area
is constant since the map projection underlying the CoE reference grid system
is an equal-area map projection).
The area estimator for a given drainage region \,$d$\, and reference land cover type \,$r$\,
is thus a weighted sum of Horvitz-Thompson mean estimators with non-random weights and
mutually independent summands; its variance estimate was then computed accordingly. 

\vskip 0.3cm
\noindent
Precisely, the area \,$\mathcal{A}^{(r)}_{d}$\, of reference land cover \,$r$\,
in drainage region \,$d$\, was estimated as follows:
\vskip -0.3cm

%\renewcommand*{\ref}{\textnormal{ref}}
%\newcommand*{\map}{\textnormal{map}}

%%%%%%%%%%

\begin{eqnarray*}
\mathcal{A}^{(r)}_{d}
	& = &
	\underset{p}{\sum}\;
	\mathcal{A}^{(r)}_{p\,\cap\,d}
\;\; = \;\;
	\underset{p}{\sum}\;
	\overset{K}{\underset{k=1}{\sum}}\;\,
	\mathcal{A}^{(r)}_{(p,k)\,\cap\,d}
\\
	& = &
	{\color{blue}
	\underset{p}{\sum}\;
	\overset{K}{\underset{k=1}{\sum}}\;\,
	\mathcal{A}_{(p,k)\,\cap\,d}
	%A_{\textnormal{{\tiny cell}}}
	%\,\cdot\,
	%M_{r,j}
	\cdot
	P\!\left(
		\left. C_{\textnormal{ref}} \overset{{\color{white}1}}{=} r \,\;\right\vert\!\!
		\begin{array}{c} R = p\,\cap\,d, \\ C_{\textnormal{map}} = k \end{array}
		\!\!\right)}
% & {\color{red}\approx} &
\;\; {\color{red}\approx} \;\;
	{\color{blue}
	\underset{p}{\sum}\;
	\overset{K}{\underset{k=1}{\sum}}\;\,
	\mathcal{A}_{(p,k)\,\cap\,d}
	%A_{\textnormal{{\tiny cell}}}
	%\,\cdot\,
	%M_{r,j}
	\cdot
	{\color{red}\widehat{P}}\!\left(
		\left. C_{\textnormal{ref}} \overset{{\color{white}1}}{=} r \,\;\right\vert\!\!
		\begin{array}{c} {\color{red}R = p}, \\ C_{\textnormal{map}} = k \end{array}
		\!\!\right)}
\\
& = &
	\underset{p}{\sum}\;
	\overset{K}{\underset{k=1}{\sum}}\;\,
	\mathcal{A}_{(p,k)\,\cap\,d}
	%A_{\textnormal{{\tiny cell}}}
	%\,\cdot\,
	%M_{r,j}
	\,\cdot
	\dfrac{
		\left(\,\begin{array}{c}
			%\textnormal{\tiny \# grid cells of reference class \,$k$}
			\textnormal{\small land area of reference class \,$r$}
			\\
			\textnormal{\small in selected grid cells}
			\\
			\textnormal{\small of map class \,$k$\, in ecoprovince \,$p$}
			\end{array}\,\right)
		}{
		\left(\,\begin{array}{c}
			%\textnormal{\tiny \# grid cells}
			\textnormal{\small total area (land or otherwise)}
			\\
			\textnormal{\small of selected grid cells}
			\\
			\textnormal{\small of map class \,$k$\, in ecoprovince \,$p$}
			\end{array}\,\right)
		}
% \\
% & = &
\;\; = \;\;
	\underset{p}{\sum}\;\,
	\overset{K}{\underset{k=1}{\sum}}\;\,
	\mathcal{A}_{(p,k)\,\cap\,d}
	%A_{\textnormal{{\tiny cell}}}
	%\,\cdot\,
	%M_{r,j}
	\,\cdot
	\left(\,
		\dfrac{
			\underset{i\,\in\mathcal{S}_{p,k}}{\sum} y^{(r)}_{i}
			}{
			{\color{white}.}
			A_{\textnormal{{\tiny cell}}}
			\cdot
			\left\vert\;
				\overset{{\color{white}.}}{\mathcal{S}_{p,k}}
				\;\right\vert
			{\color{white}.}
			}
		\,\right)
\\
& = &
	\underset{p}{\sum}\;\,
	\overset{K}{\underset{k=1}{\sum}}\;
	\left(\,
		\dfrac{
			\mathcal{A}_{(p,k)\,\cap\,d} %A^{(p,k)}_{d}
			}{
			\mathcal{A}_{\textnormal{{\tiny cell}}}
			}
		\,\right)
	\,\cdot
	\left(\;
		\dfrac{1}{\vert\,\overset{{\color{white}.}}{\mathcal{S}_{p,k}}\,\vert}
		\cdot\!
		\underset{i\,\in\mathcal{S}_{p,k}}{\sum}\,y^{(r)}_{i}
		\,\right)
% \\
% & = &
\;\; = \;\;\;
	\underset{p}{\sum}\;\,
	\overset{K}{\underset{k=1}{\sum}}\;
	\left(\,
		\dfrac{
			\mathcal{A}_{(p,k)\,\cap\,d} %A^{(p,k)}_{d}
			}{
			\mathcal{A}_{\textnormal{{\tiny cell}}}
			}
		\,\right)
	\,\cdot\,
	\widehat{\overline{y}}^{(r)}_{p,k}
	% \left(\;
	% 	\dfrac{1}{\vert\,\overset{{\color{white}.}}{\mathcal{S}_{p,k}}\,\vert}
	% 	\cdot\!
	% 	\underset{i\,\in\mathcal{S}_{p,k}}{\sum}\,y^{(r)}_{i}
	% 	\,\right)
\end{eqnarray*}
\noindent
where
\begin{equation*}
\widehat{\overline{y}}^{(r)}_{p,k}
\;\;\; := \;\;\;
	\dfrac{1}{\vert\,\overset{{\color{white}.}}{\mathcal{S}_{p,k}}\,\vert}
	\cdot\!
	\underset{i\,\in\mathcal{S}_{p,k}}{\sum}\,y^{(r)}_{i}
\;\;\; = \;\;\;
	\left\{\begin{array}{c}
	\textnormal{\small Horvitz-Thompson estimate of}
	\\
	\textnormal{\small the average land area of reference class $r$}
	\\
	\textnormal{\small for grid cells of map class $k$ in ecoprovince $p$}
	\end{array}\right\}
\end{equation*}
\begin{multicols}{2}

	\begin{center}
	\begin{minipage}{3.25in}
	\begin{eqnarray*}
	\mathcal{A}^{(r)}_{d}
	& := &
		\left\{\begin{array}{c}
		\textnormal{\small land area}
		\\
		\textnormal{\small of reference class $r$}
		\\
		\textnormal{\small in drainage region $d$}
		\end{array}\right\}
	\\ \\ \\
	\mathcal{A}^{(r)}_{p\,\cap\,d}
	& := &
		\left\{\begin{array}{c}
		\textnormal{\small land area}
		\\
		\textnormal{\small of reference class $r$}
		\\
		\textnormal{\small in the intersection of}
		\\
		\textnormal{\small ecoprovince $p$}
		\\
		\textnormal{\small and drainage region $d$}
		\end{array}\right\}
	\end{eqnarray*}
	\end{minipage}
	\end{center}

\columnbreak

	\begin{center}
	\begin{minipage}{3.25in}
	\begin{eqnarray*}
	\mathcal{A}^{(r)}_{(p,k)\,\cap\,d}
	& := &
		\left\{\begin{array}{c}
		\textnormal{\small land area}
		\\
		\textnormal{\small of reference class $r$ and map class $k$}
		\\
		\textnormal{\small in the intersection}
		\\
		\textnormal{\small of ecoprovince $p$ and drainage region $d$}
		\end{array}\right\}
	\\
	\mathcal{A}_{(p,k)\,\cap\,d}
	& := &
		\left\{\begin{array}{c}
		\textnormal{\small land area of map class $k$}
		\\
		\textnormal{\small in the intersection}
		\\
		\textnormal{\small of ecoprovince $p$ and drainage region $d$}
		\end{array}\right\}
	\\
	\mathcal{A}_{\textnormal{cell}}
	& := &
		\left\{\begin{array}{c}
		\textnormal{\small (constant) area}
		\\
		\textnormal{\small of a grid cell}
		\end{array}\right\}
	\end{eqnarray*}
	\end{minipage}
	\end{center}

\end{multicols}
\vskip -0.5cm
\begin{eqnarray*}
P\!\left(
	\left.
		C_{\textnormal{ref}} \overset{{\color{white}1}}{=} r
		\,\;\right\vert\!
	\begin{array}{c} R = p \cap d \\ C_{\textnormal{map}} = k \end{array}
	\,\right)
& := &
	\left\{\begin{array}{c}
	\textnormal{\small conditional probability}
	\\
	\textnormal{\small for an arbitrary point in a grid cell}
	\\
	\textnormal{\small to be of reference class $r$}
	\\
	\textnormal{\small given that the grid cell lies in the intersection of}
	\\
	\textnormal{\small ecoprovince $p$ and drainage region $d$}
	\\
	\textnormal{\small and has map class $k$}
	\end{array}\right\}
\\
\widehat{P}\!\left(
	\left.
		C_{\textnormal{ref}} \overset{{\color{white}1}}{=} r
		\,\;\right\vert\!
	\begin{array}{c} R = p \\ C_{\textnormal{map}} = k \end{array}
	\,\right)
& := &
	\overset{{\color{white}\textnormal{\large$1$}}}{
		\left\{\begin{array}{c}
		\textnormal{\small estimated conditional probability}
		\\
		\textnormal{\small for an arbitrary point in a grid cell}
		\\
		\textnormal{\small to be of reference class $r$}
		\\
		\textnormal{\small given that the grid cell lies in ecoprovince $p$}
		\\
		\textnormal{\small and has map class $k$}
		\end{array}\right\}
	}
\end{eqnarray*}

% 	\begin{equation*}
% 	d = \textnormal{DR or SDA},
% 	\;
% 	k = \textnormal{reference class},
% 	\end{equation*}
% 	\begin{equation*}
% 	{\color{white}22222222.}
% 	r = \textnormal{ecoprovince},
% 	\;
% 	j = \textnormal{map class},	
% 	\;
% 	(p,k) = \textnormal{stratum},
% 	\end{equation*}
% %	\begin{eqnarray*}
% %	M_{r,j}
% %	& := &
% %		\left\{\begin{array}{c}
% %		\textnormal{\tiny number of grid cells}
% %		\\
% %		\textnormal{\tiny of map class $j$}
% %		\\
% %		\textnormal{\tiny in ecoprovince $r$}
% %		\end{array}\right\}
% %	\end{eqnarray*}
% 	\vskip -0.25cm
% 	\begin{eqnarray*}
% 	\widehat{P}\!\left(
% 		\left. C_{\textnormal{ref}} \overset{{\color{white}1}}{=} k \,\;\right\vert\, R = r, C_{\textnormal{map}} = j
% 		\,\right)
% 	\\
% 	& := &
% 		\left\{\begin{array}{c}
% 		\textnormal{\tiny estimated conditional probability}
% 		\\
% 		\textnormal{\tiny for an arbitrary point in a grid cell}
% 		\\
% 		\textnormal{\tiny to be of reference class $k$}
% 		\\
% 		\textnormal{\tiny given that the grid cell lies in ecoprovince $r$}
% 		\\
% 		\textnormal{\tiny and has map class $j$}
% 		\end{array}\right\}
% 	\end{eqnarray*}
% 	\vskip -0.5cm
% 	\begin{eqnarray*}
% 	A^{(\cdots)}_{\;\,\cdots}
% 	\!\!\!\!
% 	& := &
% 		\left\{\,\begin{array}{c}
% 		\textnormal{\tiny area of}
% 		\\
% 		\textnormal{\tiny $\cdots$}
% 		\end{array}\,\right\}
% 	\\
% 	\overset{{\color{white}\textnormal{\Large1}}}{y^{(r)}_{i}}
% 	{\color{white}22}
% 	\!\!\!\!
% 	& := &
% 		\left\{\begin{array}{c}
% 		\textnormal{\tiny land area in the $i^{\textnormal{\,th}}$ grid cell}
% 		\\
% 		\textnormal{\tiny of reference class $k$}
% 		\end{array}\!\right\}
% 	\\
% 	\overset{{\color{white}\textnormal{\Large1}}}{\mathcal{S}_{r,j}}
% 	{\color{white}22}
% 	\!\!\!\!
% 	& := &
% 		\left\{\,\begin{array}{c}
% 		\textnormal{\tiny sample of grid cells}
% 		\\
% 		\textnormal{\tiny selected from \,$\mathcal{U}_{r,j}$}
% 		\end{array}\,\right\}
% 	\end{eqnarray*}

%%%%%%%%%%


          %%%%% ~~~~~~~~~~~~~~~~~~~~ %%%%%

\subsection{Area estimates for sub-drainage areas $\times$ reference land cover types and for major 
drainage areas $\times$ reference land cover types}
This aforementioned area estimation approach for drainage regions $\times$
reference land cover types is also directly applicable to sub-drainage area $\times$
reference land cover types, as well as for major drainage areas $\times$ reference land cover types.
We applied the same estimation procedure to generate these two sets of area estimates. 

          %%%%% ~~~~~~~~~~~~~~~~~~~~ %%%%%

\subsection{Calibration and raking} 
There were two CODR tables in the 2025/03/27 release:
\renewcommand{\theenumi}{(\roman{enumi})}
\renewcommand{\labelenumi}{\mbox{}\;\theenumi$\;$}
\begin{enumerate}
\item
    area estimates by reference land cover type and ecoprovince/ecozone,
\item
    area estimates by reference land cover type and drainage region/sub-drainage area/major drainage area.
\end{enumerate}
Note that:
\begin{itemize}
\item
    both tables contain land-cover-specific area estimates at the Canada level,
\item
    the areas of the target geographies (ecoprovinces, drainage regions, etc.) are known \textit{a priori}.
\end{itemize}

\vskip 0.3cm
\noindent
Calibration was performed on Table (i)
to ensure, for each given ecoprovince,
the area estimates of different reference land cover types within that ecoprovince
sum to the \textit{a priori} known area of that ecoprovince.

\vskip 0.5cm
\noindent
Two-way raking was performed on Table (ii) to ensure that 
\begin{itemize}
\item
    the Canada-level land-cover-specific area estimates in Table (ii) agree with those in Table (i), and
\item
    for each given drainage region/sub-drainage area,
    the area estimates of different reference land cover types within that drainage region/sub-drainage area
    sum to the \textit{a priori} known area of that drainage region/sub-drainage area.
\end{itemize}

          %%%%% ~~~~~~~~~~~~~~~~~~~~ %%%%%
